\begin{abstract}
Automated bearing fault diagnosis relies on trustworthy predictive confidence to determine when the system should
defer to a human operator. Machine-learning classifiers are usually trained and calibrated on artificial defects but
must ultimately diagnose naturally developing damage, such as fatigue pitting. Whether their confidence can still be
trusted after this shift has received little attention. This study evaluates whether confidence thresholds, post-hoc
calibration and conformal prediction guarantees transfer from artificial to real damage. Random Forest, support vector
machine and XGBoost classifiers on envelope-spectrum features and a one-dimensional convolutional neural network
(1D-CNN) on raw vibration are trained and calibrated on artificially damaged Paderborn bearings and tested on real
damage, with bearing-level splits, 105 choices of calibration bearings and four operating conditions. At the main
operating condition, the feature-based models are close to calibrated on unseen artificial bearings (expected
calibration error 0.01--0.09), but on real damage accuracy drops from 0.88--0.91 to about 0.69, and thresholds tuned
for a 5\,\% target error give 24--27\,\% error, exceeding the target for every calibration choice. Neither temperature
scaling nor isotonic regression corrects this overconfidence, and conformal coverage falls to 0.65--0.68 against a
nominal 0.90. The failures stem largely from ``silent'' real-damage bearings lacking distinct fault-frequency peaks in
their envelope spectra, which the models confidently call healthy; they recur at the other operating conditions, more
mildly where the features are weaker. Calibrating on two labelled real bearings per class reduces the error to
0.10--0.11 but meets the target in only about half of the trials and automates far fewer cases. The 1D-CNN fails even
on unseen artificial bearings (0.46 accuracy vs.\ 0.98 under random window splits), and adaptive batch normalisation
(AdaBN) on unlabelled real-damage data raises its real-damage accuracy from 0.43 to 0.51 but not the reliability of
its confidence. Consequently, safety-critical automation thresholds must be evaluated on real-damage data with strict
bearing-level validation, and results reported per bearing.
\end{abstract}
